\scanpage{18}
\mspage{9}
\uline{Définition de $\alpha$}\qquad On va définir
\[
\alpha(\unsure{\mathcal E|\Sigma},\mathcal G,\pi,\mathfrak Z\ldots)=(\Pi,\widetilde\Pi,p,\mathcal G,\unsure{a|\sim})
\]
\[
\left\{\begin{array}{l}
\operatorname{Ob}\Pi=\mathfrak Z/\pi\unsure{.}\quad(\text{soit }J)\\
\operatorname{Isom}(j,j')=\operatorname{Isom}_{\pi}(\mathfrak Z_j,\mathfrak Z_{j'})\\
\text{composition des flèches évidente}
\end{array}\right.\qquad\bigg|\quad\text{d'où groupoïde }\Pi
\]
\[
\deleted{\widetilde\Pi_j=\mathfrak Z_j}\qquad\operatorname{Ob}\widetilde\Pi=\mathfrak Z
\]
\diagram{\begin{tikzcd}[ampersand replacement=\&]
\operatorname{Isom}(z,z')=\{\mathrm{pt}\}\arrow[d]\&\begin{gathered}\forall z,z'\in\mathfrak Z\\(\widetilde\Pi\text{ cat. discrète}\ldots)\end{gathered}\\
\operatorname{Isom}(j,j')\simeq\operatorname{Isom}(\mathfrak Z_j,\mathfrak Z_{j'})\&
\end{tikzcd}}
\deleted{(\unsure{où} $j,j'$ \unsure{les} \unsure{images}\unsure{)}}\inserted{(si $j=\dot z$, $j'=\dot z'$)}
\margin{\figure{\begin{tikzpicture}[x=1cm,y=1cm]
\draw (0,0.20)--(0,-2.95);
\node[anchor=west,inner sep=0pt] at (0.15,0) {$\widetilde\Pi$ cat. \uline{discrète}};
\node[anchor=west,inner sep=0pt] at (0.15,-0.38) {sur $\Pi$ définie par};
\node[anchor=west,inner sep=0pt] at (0.15,-0.86) {foncteur};
\node[anchor=west,inner sep=0pt] at (0.95,-0.60) {\inserted{\unsure{contr.}}};
\draw[thin] (0.94,-0.64) .. controls (0.75,-0.66) and (0.57,-0.74) .. (0.45,-0.83);
\node[inner sep=0pt] at (0,-1.35) {$\Longleftarrow$};
\node[anchor=west,inner sep=0pt] at (0.15,-1.45) {\begin{tikzcd}[ampersand replacement=\&,column sep=small]
F:\Pi^\circ\arrow[r]\&\operatorname{Ens}
\end{tikzcd}};
\node[anchor=west,inner sep=0pt] at (0.15,-2.10) {$F(j)=\mathfrak Z_j,\ldots$};
\node[anchor=west,inner sep=0pt] at (0.15,-2.72) {$\widetilde\Pi\simeq\Pi_{/F}$};
\end{tikzpicture}}}
$\downarrow$ transforme l'unique \unsure{pt|p\textsuperscript{t}} de $\operatorname{Isom}(z,z')$ en l'unique \inserted{$\pi$-}isom $\mathfrak Z_j\mathbin{-}\mathfrak Z_{j'}$ transformant $z$ en $z'\ldots.$

Pour l'instant on n'a pas utilisé $\unsure{\mathcal E|\Sigma}$ et l'op. de $\unsure{\mathcal E|\Sigma}$ sur $\mathfrak Z$. On va l'utiliser maintenant pour faire opérer $\mathcal G=\unsure{\mathcal E|\Sigma}/\pi$ sur $\Pi$, et $\unsure{\mathcal E|\Sigma}$ sur $\widetilde\Pi$, de façon compatible.
Comme $\unsure{\mathcal E|\Sigma}$ opère sur $\mathfrak Z=\operatorname{Ens}(\widetilde\Pi)$, il opère en fait sur $\widetilde\Pi$, cat. discrète associée, \deleted{\unsure{opérant librement}} et \deleted{le} $\pi$ opère par $\Pi$-auto---, d'où l'opération de $\mathcal G$ sur $\pi$ par passage au quotient.
\mspage{17}
